\documentclass{article}

% AI4Mat (NeurIPS workshop) typically uses the NeurIPS style file in
% "final" or workshop mode. Drop the official neurips_2026.sty here and
% adjust the option below when the CFP is out.
\usepackage[preprint]{neurips_2026}

\usepackage[utf8]{inputenc}
\usepackage[T1]{fontenc}
\usepackage{hyperref}
\usepackage{url}
\usepackage{booktabs}
\usepackage{amsfonts,amsmath,amssymb}
\usepackage{nicefrac}
\usepackage{microtype}
\usepackage{graphicx}
\usepackage{xcolor}

\newcommand{\todo}[1]{\textcolor{red}{[TODO: #1]}}

\title{Structure-to-Recipe: Calibrated Inverse Composition-Process Prediction
for Mg Alloys}

\author{%
  \todo{Author list and affiliations}%
}

\begin{document}

\maketitle

\begin{abstract}
Closing the inverse loop of integrated computational materials engineering
requires answering the question a forward simulation cannot: which chemistry
and processing route produced an observed microstructure. We study this
structure-to-recipe problem on 108 as-extruded conditions spanning 14
magnesium alloys and make three findings. First, on a frozen
alloy-stratified random split over whole conditions, the two halves of a
recipe favour different representations: conventional statistical descriptors
identify the alloy best, whereas graph embeddings of the grain structure
predict the extrusion window best, are the most stable across heads, and give
the only near-calibrated process uncertainty. Second, the literature-default
continuous output is the wrong type for chemistry: alloy compositions live on
a small discrete lattice, and a probabilistic 14-way classifier fused with
retrieval and condition balancing cuts the present-element WAPE of direct
continuous regression from 40.7\% to 24.1\% at equal element MAE, while
additionally returning alloy identity and a distribution over nominal
compositions as the natural, calibratable uncertainty statement. Third, we
show why the split unit is decisive in this repeated-image dataset: random
image splitting gives 100\% alloy accuracy and zero composition error because
every test condition is also represented in training. We therefore estimate
generalization by assigning whole conditions to the frozen split.
\end{abstract}

\section{Introduction}
\label{sec:intro}
% Story arc (strategy/00): forward ICME is mature; the inverse direction,
% inferring provenance from structure, is the bottleneck for alloy discovery,
% failure analysis, and closing design loops.
\todo{Paragraph 1: forward vs inverse ICME, why provenance inference matters.}
\todo{Paragraph 2: why point regression is insufficient; the map is one-to-many,
labels are jointly dependent (composition co-determines viable extrusion
windows); need joint, calibrated densities.}
\todo{Paragraph 3: contributions list.}

\textbf{Contributions.}
(i) A unified benchmark: three descriptor families, identical heads, metrics
and splits on 108 extrusion conditions of 14 Mg alloys, over the two halves of
a recipe, namely composition from structure and the extrusion window given the
identified composition.
(ii) A systematic head study showing that discrete composition prediction
(condition-balanced classification over the composition lattice fused with
retrieval) dominates direct continuous regression for recipe recovery in the
scarce-data regime, with the fused class distribution as the composition
uncertainty statement.
(iii) A split-unit audit showing that image-random evaluation is a memorization
upper bound, and a leakage-safe benchmark that assigns whole processing
conditions to a frozen alloy-stratified split.

\section{Related work}
\label{sec:related}
\todo{Microstructure descriptors and MKS (Kalidindi et al.); inverse design
generative models; normalizing flows for SBI/posterior estimation; GNNs on
grain graphs; prior forward work on this dataset (Acta Materialia 2025).}

\section{Dataset and task}
\label{sec:data}
% From strategy/03: 108 as-extruded conditions, 14 base alloys, labels =
% 8 element concentrations (wt.%) + extrusion temperature + ram speed.
\todo{Describe database, optical micrographs at four magnifications, XRD
texture; labels table; the two-level extrusion-ratio encoding; the frozen
alloy-stratified random condition split; leakage
discipline and the global-reduction limitation.}

\section{Method}
\label{sec:method}

\subsection{Descriptor pipelines}
\todo{Conventional: one 438-D vector per image, comprising 3-point PCA (120),
Gram Isomap (200), aspect-ratio and equivalent-diameter histograms (29 each),
and condition-level GSH Isomap features (60). Images from the 90, 100, 120,
and 150\,\textmu m fields of view are repeated observations, not concatenated
feature blocks. GenAI: frozen ViT encoder latents, 16 tokens $\times$ 80-D.
GNN: grain-adjacency graphs from DREAM.3D RVEs, encoded by GATv2 with
multi-seed attention pooling; the main comparison uses the 90\,\textmu m
graphs.}

\subsection{Autoregressive flow-transformer head}
% From strategy/01: 10 label tokens, causal decoder d=128, 4 layers,
% cross-attention to descriptor tokens, hurdle + RQ spline flows.
\todo{Factorization equation, token order, hurdle likelihood, spline flows,
tolerance dequantization, training recipe, inference (S=256 ancestral samples,
joint-MAP, k-medoids shortlist).}

\begin{equation}
p(\mathbf{y} \mid \mathbf{c}) \;=\; \prod_{k=1}^{10}
p\!\left(y_k \,\middle|\, y_{<k}, \mathbf{c}\right), \qquad
p(y_e \mid \cdot) = (1-\pi_e)\,\delta_0(y_e) + \pi_e\, f_e(y_e \mid \cdot).
\end{equation}

\subsection{Condition-balanced discrete composition head}
The proposed composition head treats the 14 nominal alloy compositions as a
discrete lattice. During CatBoost training, image $i$ from condition $c$
receives weight $1/n_c$, normalized to unit mean, so every processing
condition contributes equal total loss despite having 15--108 images. A
second posterior is obtained from cosine $k$-nearest-neighbor retrieval in
training-standardized descriptor space, with $k=5$, similarity temperature
0.05, and the same condition-balanced vote weights. Their image-level class
posteriors are fused as
\begin{equation}
  p_{\mathrm{fuse}}(a\mid x) = \tfrac{1}{2}p_{\mathrm{CB}}(a\mid x)
  + \tfrac{1}{2}p_{\mathrm{kNN}}(a\mid x),
\end{equation}
To reduce long-tail bias, let $N_a$ be the number of unique training
conditions for alloy $a$. Before condition-level pooling, we remove a
validation-selected power of this empirical prior from each image posterior,
\begin{equation}
  p_{\tau}(a\mid x) =
  \frac{p_{\mathrm{fuse}}(a\mid x) / N_a^{\tau}}
       {\sum_b p_{\mathrm{fuse}}(b\mid x) / N_b^{\tau}},
  \qquad \tau\in\{0,0.25,0.5,0.75,1\}.
\end{equation}
The corrected probabilities are then averaged over all images of the query
condition. The exponent $\tau$ is selected on the validation conditions only,
by WAPE with MAE as the tie-breaker, and the selected rule is evaluated once
on the untouched test conditions; Table~\ref{tab:heads} shows that this
selection is not reliable at the present validation size. The point recipe is
the nominal composition at
$\arg\max_a p_{\tau}(a\mid x)$; the full 14-way distribution is retained for
uncertainty analysis. The model receives only microstructure and texture
descriptors as inputs.

\section{Experiments}
\label{sec:experiments}

\subsection{Setup}
\todo{Metrics: present-element WAPE, element MAE, class NLL, top-$k$ alloy
identification for composition; MAE, WAPE, $R^2$ and 90\% coverage for the
extrusion window. Present-only WAPE = sum of absolute errors over sum of true
values; pointwise MAPE is unusable because trace-level targets explode the
ratio.}
All experiments use one frozen alloy-stratified random split of the 108 whole
conditions (seed 0) into 72 training, 18 validation and 18 test conditions
(2{,}909/729/710 images), with no condition shared between parts. All heads
see identical splits and identical per-condition aggregation. Composition
heads receive the representation and a two-level encoding of the extrusion
ratio; process heads additionally receive the known element weight fractions.
Because the estimate is a single split over 18 test conditions, top-$k$ moves
in steps of $1/18$ and small head-to-head differences are not resolvable; we
report them for completeness but claim only the effects that survive this
resolution.

\subsection{The head study: recipes live on a lattice}
\begin{table}[t]
  \caption{Head study on conventional descriptors, frozen seed-0 split.
  The last two rows regress the eight element fractions directly and cannot
  name an alloy. Treating composition as a discrete lattice wins on WAPE and
  on identification, while element MAE is essentially tied with continuous
  regression.}
  \label{tab:heads}
  \centering
  \begin{tabular}{lccccc}
    \toprule
    Head & El.\ WAPE $\downarrow$ & El.\ MAE $\downarrow$ &
    Top-1 $\uparrow$ & Top-3 $\uparrow$ & Class NLL $\downarrow$ \\
    \midrule
    kNN retrieval ($k{=}5$)  & 31.3\% & 0.433 & 0.39 & 0.50 & 6.04 \\
    XGBoost classifier       & 51.6\% & 0.392 & 0.33 & 0.56 & 3.87 \\
    CatBoost classifier      & 24.8\% & 0.284 & 0.44 & 0.61 & 2.40 \\
    FT-Transformer           & 42.2\% & 0.266 & 0.44 &
    $\mathbf{0.67}$ & 2.31 \\
    CatBoost{+}kNN fusion    & 25.3\% & 0.307 & 0.44 & 0.56 & 2.29 \\
    Condition-balanced fusion & $\mathbf{24.1\%}$ & 0.248 &
    $\mathbf{0.50}$ & 0.56 & 2.14 \\
    Prior-adjusted balanced fusion & 47.7\% & 0.297 & 0.39 & 0.56 &
    $\mathbf{2.07}$ \\
    XGBoost regression       & 40.7\% & $\mathbf{0.240}$ & -- & -- & -- \\
    CatBoost regression      & 45.7\% & 0.290 & -- & -- & -- \\
    \bottomrule
  \end{tabular}
\end{table}
The eight element fractions are not free coordinates: they take only 14 joint
values, so a head that scores those values directly is solving an easier and
better-posed problem than one that regresses each fraction. On the frozen
split the condition-balanced fusion reaches 24.1\% present-element WAPE
against 40.7\% for the best direct regression, and it additionally names the
alloy, which the regressors cannot. Element MAE does not separate the two
families (0.248 versus 0.240): the gain is concentrated in the elements that
are actually present, which is what WAPE measures and what a recipe needs.
The retained 14-way distribution is the composition uncertainty statement
(class NLL 2.07--2.14 for the fused heads versus 6.04 for retrieval alone).
Two negative results are worth recording. The validation-selected prior
correction, which helped under repeated resampling, does not survive on a
single split: it degrades WAPE to 47.7\% while still giving the best NLL, so
$\tau$ cannot be selected reliably from 18 validation conditions. The
FT-Transformer attains the best top-3 and a competitive MAE but a poor WAPE,
consistent with it spreading mass over chemically similar alloys.
\todo{Add conformal recalibration of the fused distribution.}

\subsection{Which representation carries provenance information?}
Predicting a recipe has two halves that are not symmetric. The alloy must be
identified from structure alone, whereas the extrusion parameters are only
useful once the material is known, so the process heads are conditioned on the
known element weight fractions in addition to the representation. Both halves
also receive a two-level encoding of the extrusion ratio, which separates the
Mg--Gd series from the remaining alloys. Table~\ref{tab:cross} evaluates the
same head twice across all three pipelines on the frozen seed-0 split: the
condition-balanced fusion for composition and the exact GP for the process
window.

\begin{table}[t]
  \caption{Cross-pipeline comparison on the frozen alloy-stratified random
  split (seed 0; 72/18/18 conditions). Composition uses the
  condition-balanced CatBoost{+}kNN fusion; the process window uses the exact
  GP and is conditioned on the known composition, so $T$ errors are not
  comparable to Table~\ref{tab:heads}. Single split over 18 test conditions.}
  \label{tab:cross}
  \centering
  \begin{tabular}{lcccccc}
    \toprule
    & \multicolumn{3}{c}{Composition} &
    \multicolumn{3}{c}{Process $\mid$ composition} \\
    \cmidrule(lr){2-4}\cmidrule(lr){5-7}
    Representation & El.\ MAE $\downarrow$ & Top-1 $\uparrow$ &
    El.\ WAPE $\downarrow$ &
    $T$ MAE $\downarrow$ & $T$ $R^2$ $\uparrow$ & cov$_{90}$ \\
    \midrule
    Conventional & $\mathbf{0.248}$ & $\mathbf{0.50}$ & $\mathbf{24.1\%}$ &
    42.9 & 0.230 & 0.78 \\
    GenAI latents & 0.266 & 0.39 & 33.5\% & 43.5 & 0.113 & 0.72 \\
    GNN (90\,\textmu m) & 0.403 & 0.22 & 59.5\% &
    $\mathbf{32.0}$ & $\mathbf{0.524}$ & $\mathbf{0.94}$ \\
    \bottomrule
  \end{tabular}
\end{table}

The two halves rank the representations differently. Conventional descriptors
are the strongest composition representation under the head that won
Table~\ref{tab:heads}. For the process window the GNN
embeddings are both the most accurate and the only near-calibrated option, and
they are markedly more stable across heads: every GNN head lands between 0.315
and 0.524 in $T$ $R^2$ and between 0.498 and 0.575 in $v$ $R^2$, the tightest
band of the three representations, while the conventional descriptors spread
from 0.225 to 0.483 in $T$ $R^2$ and the GenAI latents fall to $-0.001$. A
graph over grains appears to retain the
deformation-rate information that the aggregated descriptor histograms
average away. With 18 test conditions the ordering of heads within a pipeline
is inside the noise, so only the cross-pipeline pattern is claimed.

The extrusion speed is learned in log space: its values follow a near
geometric grid spanning $1.18$ decades against $0.40$ for the temperature, so a
squared error on the raw scale shrinks the slow conditions toward the mean.
The GP therefore treats $v$ as lognormal, which gives a multiplicative
$90\%$ interval and $90\%$ coverage of $0.94$ on both targets for the GNN,
against $0.72$--$0.78$ on $T$ for the other two representations.

\subsection{Why the split unit matters}
Random splitting of the 4,348 individual images produces 100\%
alloy top-1 accuracy and zero composition WAPE with the same
CatBoost{+}kNN head. This is not evidence of generalization: every test
condition also occurs in training, so composition and process labels can be
recovered by matching repeated observations. All new headline experiments
therefore use one frozen alloy-stratified random split over the 108 whole
conditions (seed 0), with no condition shared across train, validation, and
test.

\subsection{Ablations and analysis}
\todo{Calibration reliability of the fused classifier; descriptor-block
importance; posterior visualizations for ambiguous conditions. Full head and
decoder ablation in the appendix (results/archive).}

\section{Discussion and limitations}
\label{sec:discussion}
\todo{Small-data regime; single alloy system (Mg, extrusion); descriptors from
2D sections; synthetic-vs-experimental graph domain gap; outlook to closing the
full design loop with the forward models.}

\section{Conclusion}
\label{sec:conclusion}
\todo{One paragraph.}

\begin{ack}
\todo{Funding and compute acknowledgements.}
\end{ack}

\bibliographystyle{plainnat}
\bibliography{references}

\appendix

\section{Architecture and training details}
\todo{Full head hyperparameters; parameter counts; standardization details.
The FT-Transformer emits one token per input feature, so the 1{,}282-D
generative latents are trained with gradient accumulation to bound attention
memory at a fixed effective batch of 256.}

\section{Additional results}
\todo{Per-label metric tables (A1-A7 from strategy/02); the full head grid for
all three pipelines (Table 3 of results/paper\_results.md); sensitivity of the
conclusions to the split seed.}

\end{document}
